\documentclass[../../main.tex]{subfiles}
\begin{document}

{\bf [解答]}

$k$ ($1 \le k \le N-1$) 回目に大吉を引く，つまり$k$回目で引くのをやめる確率$p_k$は
\begin{align}
  p_k = (1-p)^{k-1}p \quad (1 \le k \le N-1) \label{eq:1}
\end{align}
であり，$N$回目のおみくじを引く確率は$N-1$回目まで大吉を引かない時だから
\begin{align}
  p_N = (1-p)^{N-1} \label{eq:2}
\end{align}
である．従って\cref{eq:1,eq:2}より，おみくじを引く回数の期待値は
\begin{align}
  E
  &= \sum_{k=1}^N k p_k \\
  &= p\sum_{k=1}^{N-1} k(1-p)^{k-1} + N(1-p)^{N-1} \label{eq:3}
\end{align}
である．ここで
\begin{align}
  S = \sum_{k=1}^{N-1} k(1-p)^{k-1}
\end{align}
とおくと
\begin{align}
  \begin{cases}
    S  = 1 + 2(1-p) + \cdots + (N-1)(1-p)^{N-2} \\
    (1-p)S = (1-p) + \cdots + (N-2)(1-p)^{N-2} + (N-1)(1-p)^{N-1}
  \end{cases}
\end{align}
だから，辺々引いて
\begin{align}
  pS
  &= \sum_{k=0}^{N-2} (1-p)^k - (N-1)(1-p)^{N-1} \\
  &= \dfrac{1-(1-p)^{N-1}}{1-(1-p)} - (N-1)(1-p)^{N-1} \\
  &= \dfrac{1-(1-p)^{N-1}}{p} - (N-1)(1-p)^{N-1}
\end{align}
となる．これを\cref{eq:3}に代入して
\begin{align}
  E
  &= \dfrac{1-(1-p)^{N-1}}{p} - (N-1)(1-p)^{N-1} + N(1-p)^{N-1} \\
  &= \dfrac{1-(1-p)^{N-1}}{p} + (1-p)^{N-1} \\
  &= \dfrac{1-(1-p)^N}{p} \quad \text{(答)}
\end{align}
である．

$(p, N) = \left(\dfrac{1}{5}, 10\right)$ の時
\begin{align}
  E &= 5 \left\{ 1 - \left(\dfrac{4}{5}\right)^{10} \right\} \\
  &= 5 \left[ 1 - \left(\dfrac{8^5}{10^5}\right)^2 \right] \\
  &\fallingdotseq 5 ( 1 - 0.107374 ) \cdots \bigstar \\
  &= 4.463 \\
  &\fallingdotseq 4.46 \quad \text{(答)}
\end{align}
である．

\end{document}
